% =====================================================================
%  Chapter 11 — EARLY WORK IV: inheritance
% =====================================================================
\chapter{Early Work IV: What the Later Programme Inherited}
\label{ch:inheritance}

\begin{record}
\textbf{Purpose of this chapter.} The early work of
Chapters~\ref{ch:claim-a}--\ref{ch:audit} is retained and read forward. This
chapter says exactly what the sessions \Sn{32}--\Sn{47} took from it: which
objects were reused, which were rebuilt, which were abandoned, and --- the
question that matters most --- which parts of the early work are still
load-bearing today.
\end{record}

There are three inheritances. The first is a framework: the coherence
quantity \(\sigma\) and the algebraic identity underneath it became the
localised \(\sstar\) on which every session from \Sn{32} to \Sn{45} is built.
The second is an obstruction: a single clause in the audit's eighth defect
became the Calder\'on--Zygmund logarithm \(\Lfac\), which is now the
programme's central structural barrier and was rediscovered three separate
times before anybody connected it back to its source. The third is a
discipline: the withdrawal of \Sn{29}'s claim produced the verification rules
of Chapter~\ref{ch:house-rules}.

The first is the most useful. The second is the most important. The third is
the only one that would transfer unchanged to a different problem.

\section{The traceability table}
\label{sec:traceability}

\begin{center}
\footnotesize
\begin{longtable}{@{}p{0.20\textwidth} p{0.11\textwidth} p{0.24\textwidth} p{0.23\textwidth} p{0.11\textwidth}@{}}
\toprule
\textbf{Object (label, line)} & \textbf{Early status} &
\textbf{What happened to it} & \textbf{Where it reappears, \Sn{32}--\Sn{47}} &
\textbf{Status now} \\
\midrule
\endfirsthead
\toprule
\textbf{Object (label, line)} & \textbf{Early status} &
\textbf{What happened to it} & \textbf{Where it reappears} &
\textbf{Status now} \\
\midrule
\endhead
\bottomrule
\endfoot

\texttt{thm:routeC} (508) \newline Route~C bootstrap
  & Proved, \Sn{17}
  & Untouched by the audit; named survivor \#1
  & Still the terminal step of any closure: any \(\delta_0>0\) gives
    smoothness. Never invoked in \Sn{32}--\Sn{47} because no session produced
    a \(\delta_0\)
  & \PROVED \\[3pt]

\texttt{thm:shear-claimA} (1722) \newline Claim~A for shear data
  & Proved, \Sn{20}
  & Untouched; named survivor \#2
  & Not used; it is a solved case, not a tool
  & \PROVED \\[3pt]

\texttt{thm:near-shear} (1961)
  & Proved, \Sn{21}
  & Untouched; named survivor \#2
  & Not used
  & \PROVED \\[3pt]

\texttt{prop:lambda-zero} (1779) \newline \(\lambda_{\max}=0\) for shear
  & Proved, \Sn{20}
  & Untouched; named survivor \#3
  & Not used
  & \PROVED \\[3pt]

\texttt{lem:grad-omega-split} (4564) \newline
  \(\abs{\nabla\omega}^2=\abs{\nabla m}^2+m^2\abs{\nabla\ehat}^2\)
  & Proved, \Sn{30}
  & Named survivor \#4; ``sound foundation for a future quantitative alignment
    program''
  & \textbf{Used constantly.} It is the pointwise device in the \Sn{34}
    magnitude Caccioppoli (\texttt{eq:mag-caccioppoli}, 5836) that sends every
    transition-annulus term to the enstrophy-gradient budget with no
    absorption argument; and it is what makes the \Sn{36} Gram negativity
    (\texttt{lem:gram-negativity}, 6180) a statement about a positive
    semi-definite tensor
  & \PROVED \\[3pt]

\(\sigma=A_{\mathrm{loc}}/M^{3/2}\) \newline (\texttt{eq:sigma-def}, 4595)
  & Proved scale-invariant, \Sn{30}
  & Named survivor \#4; the exponent \(3/2\) is the unique fixed point of the
    NS scaling group
  & \textbf{Directly ancestral to \(\sstar\).}
    \texttt{prop:sigma-controls} (5424) proves \(\sstar\le16\sigma\) in two
    lines --- the surviving quantity controls the new one
  & \PROVED \\[3pt]

\texttt{thm:blowup-alignment} (4858) \newline ``Theorem E'': blow-up forces
  \(\sigma\to0\)
  & Claimed proved, \Sn{30}
  & \emph{Not} on the audit's survivors list; its proof invokes
    \texttt{thm:sigma-gronwall} (D5) and Regime~I (D6)
  & \textbf{Its statement is the programme's central open problem.} \Sn{38}
    identifies ``forcing \(\sstar(t)\to0\) near a putative Type-I
    singularity'' as the single deep gap; \Sn{39} reframes it as
    \texttt{target:disorder-}\newline\texttt{depletion} (7490). See
    \S\ref{sec:theorem-e-irony}
  & \OPEN \\[3pt]

\texttt{prop:coercive} (4611), \newline \texttt{lem:A-evol} (4664)
  & Asserted, \Sn{30}
  & Defect~D8, first half: coercivity never derived; \(\ehat\) singular on
    \(\{\omega=0\}\)
  & Never repaired, never reused. The later programme derives the direction
    equation from scratch instead (\texttt{lem:direction-eq}, 5349)
  & \OPEN \\[3pt]

\textbf{\(\norm{\nabla u}_{L^\infty}\le CM\)} \newline (D8, second half)
  & Assumed silently, \Sn{29}--\Sn{30}
  & \textbf{Identified by the audit as false}: CZ operators are unbounded on
    \(L^\infty\); the correct bound carries a logarithm
  & \textbf{Becomes \(\Lfac=\log(e+\norm{u}_{H^3}/M)\).} Cited by name in
    \texttt{rem:marginality} (5443); enters formally at
    \texttt{lem:local-}\newline\texttt{enstrophy-typeI} (5701, definition 5712); reappears at
    \texttt{prop:crude-K} (5903), \texttt{rem:s41-shortfall} (8395),
    \texttt{prop:s44-linfty-route} (9209); consolidated as Wall~1
  & \OPEN, structural \\[3pt]

\texttt{thm:sigma-gronwall} (4744)
  & Claimed, \Sn{30}
  & Defect~D5: hypothesis \(\int_0^TM\,dt<\infty\) is BKM
  & Not reused. The \Sn{32} rebuild is explicitly designed to avoid assuming
    an equivalent of its conclusion
  & \WITHDRAWN{} as a route \\[3pt]

\texttt{prop:vorticity-} \texttt{recursive-20} (4371), \newline
  \texttt{thm:uniform-} \texttt{enstrophy} (4407)
  & Claimed proved, \Sn{29}
  & Defects D1--D4
  & Never repaired. Wall~0 records that \S20 is usable only as a conditional
    small-data statement, after D2--D4 are fixed and the exponent re-derived
  & \WITHDRAWN \\[3pt]

\texttt{cor:claim-a-} \texttt{enstrophy} (4451), \newline \texttt{cor:prize-h1} (4467)
  & Claimed proved, \Sn{29}
  & Defect~D1
  & --- & \WITHDRAWN \\[3pt]

\texttt{prop:scale-recursive} (3976) \newline \S19, gradient energy
  & Proved conditional on \(u\in L^4_{\mathrm{loc}}\), \Sn{27}
  & Superseded by \S20 for having a sublinear exponent
    (\texttt{rem:prop19-gap}, 4175); \S20 then withdrawn
  & Not reused. The idea --- iterate \emph{downward} from a bound that always
    holds --- returns transformed in the \Sn{39} profile extraction, which
    also avoids needing a smallness hypothesis to start
  & \SUPERSEDED \\[3pt]

\texttt{prop:theta-cacc} (1306) \newline pressure-free Caccioppoli
  & Proved, \Sn{19}
  & Untouched by the audit
  & Not reused. Its motive --- eliminate the pressure --- is carried instead
    by \texttt{lem:direction-eq}, which is pressure-free by construction
  & \PROVED, dormant \\[3pt]

\texttt{lem:pressure-} \texttt{traceless} (2044) \newline
  \(C_{\mathrm{tr}}=\tfrac12C_{\mathrm{CZ}}\)
  & Proved, \Sn{21}
  & Untouched
  & Not reused
  & \PROVED, dormant \\[3pt]

\texttt{lem:gn-divfree} (2211), \newline \texttt{prop:mean-obstacle} (2268),
  \newline \texttt{thm:routeB-closure} (2397)
  & Proved, \Sn{22}
  & Untouched
  & Not reused. The mean-Morrey line was never returned to after \Sn{31}
  & \PROVED, dormant \\[3pt]

\texttt{thm:conditional-} \texttt{mean-morrey} (2657), \newline
  \texttt{cor:complete-} \texttt{hierarchy} (3569)
  & Proved as implications, \Sn{23}, \Sn{25}
  & Untouched; the chain is a circle and was known to be
  & Not reused
  & \PROVED, dormant \\[3pt]

\texttt{lem:pressure-} \texttt{forcing} (3672) \newline \(F_p\le Cr^{-23/15}\)
  & Proved, \Sn{26}
  & Untouched; better than the naive \(r^{-2}\), still divergent
  & Not reused
  & \PROVED, dormant \\[3pt]

The three angles of \Sn{18} \newline (CKN capacity; \(A(r)\) iteration;
  heat-kernel positivity)
  & Each obstructed, \Sn{18}
  & Untouched
  & CKN \(\varepsilon\)-regularity returns in a different role in \Sn{39}, as
    the literature that would have to supply hypothesis~(R) --- with the named
    trap that \(\varepsilon\)-regularity needs \emph{smallness} while Type-I
    supplies only \emph{boundedness}
  & \OPEN \\[3pt]

\texttt{EXP-L4-GAM-} \texttt{ADV-001} \newline \(\Gamma\) not suppressed
  & Negative, \Sn{27}
  & Correctly read as a statement about free viscous decay, not about NS
  & Prefigures Wall~6: a diagnostic whose regime does not match the
    conjecture's cannot test it
  & unresolved \\[3pt]

\texttt{EXP-L4-MM-*} \newline mean-Morrey sampling
  & \Sn{23}, \Sn{25}
  & \(47\) orders of magnitude of sampling bias, found and corrected
  & The sampling rule it produced governs \texttt{scale\_diagnostics.py} and
    \texttt{band\_concentration.py}
  & method \\[3pt]

\texttt{EXP-L4-SIGMA-*} \newline the \(\sigma\) runs
  & \Sn{30}
  & Logged to a Markdown file, not the database --- the programme's only
    breach of its own logging rule
  & Superseded by \texttt{EXP-L4-BRIDGE-*}, which measure \(\sstar\) properly
    and are in the database
  & not a result \\
\end{longtable}
\end{center}

\section{First inheritance: \texorpdfstring{\(\sigma\)}{sigma} becomes
\texorpdfstring{\(\sstar\)}{sigma-star}, and the pincer}
\label{sec:sigma-to-sstar}

The audit left exactly two sound objects on the geometric side: the split
identity and the scale invariance of \(\sigma\). \Sn{32} rebuilt on precisely
those two and nothing else.

\subsection{What was rebuilt, and what was kept}

The rebuild has four steps, and it is instructive that three of them are
\emph{re-derivations of things \Sn{30} had asserted}.

\begin{enumerate}[leftmargin=2.2em]
\item \emph{The direction equation, derived from scratch}
  (\texttt{lem:direction-eq}, line~5349). With \(\ehat=\omega/\abs{\omega}\),
  \(\alpha=\ehat\cdot S\ehat\) and \(\Dt=\partial_t+u\cdot\nabla\):
  \begin{align*}
    \Dt\abs{\omega} &= \alpha\abs{\omega}
      + \nu\bigl(\Delta\abs{\omega}-\abs{\omega}\abs{\nabla\ehat}^2\bigr),\\
    \Dt\ehat &= \nu\Delta\ehat
      + 2\nu(\nabla\log\abs{\omega}\cdot\nabla)\ehat
      + \nu\abs{\nabla\ehat}^2\ehat
      + P_{\ehat^\perp}\bigl((\ehat\cdot\nabla)u\bigr).
  \end{align*}
  The document's remark is pointed: these equations contain \emph{no
  pressure}, and this is ``the same structural advantage \S20 sought, obtained
  without any invalid estimate''. The motive of the withdrawn work was
  correct; its execution was not; the rebuild keeps the motive.
\item \emph{A correctly localised coherence quantity}
  (\texttt{def:sigma-star}, line~5393). With \(\rstar=M^{-1/2}\) and \(x^*\) a
  point where \(\abs{\omega}\ge M/2\),
  \[
    \sstar(t) = (\rstar)^{-1}\int_{B_{\rstar}(x^*)}
      \abs{\nabla\ehat}^2\,\mathbf 1_{\{\abs{\omega}\ge M/4\}}\,dx .
  \]
  \texttt{prop:sigma-star-invariant} (line~5402) proves exact scale
  invariance directly, and the calibration remark that follows fixes its
  meaning: if \(\abs{\nabla\ehat}\sim1/\rstar\) throughout the ball --- the
  borderline Constantin--Fefferman coherence at the self-similar scale --- then
  \(\sstar\sim1\). So \(\sstar\ll1\) means \emph{better than borderline}
  coherence, and \(\sstar\) is the natural small parameter for an
  \(\varepsilon\)-regularity statement.
\item \emph{The link to the survivor} (\texttt{prop:sigma-controls},
  line~5424): \(\sstar\le16\sigma\), proved in four lines. This is the formal
  moment of inheritance --- the quantity the audit certified controls the
  quantity the new programme uses.
\item \emph{The criticality observation} (\texttt{rem:marginality},
  line~5443), which is where the second inheritance enters and is discussed in
  \S\ref{sec:cz-inheritance}.
\end{enumerate}

\subsection{The pincer}

On those foundations \Sn{32} states its target, \texttt{thm:pincer-conditional}
(line~5495): given two named hypotheses --- H1, the Bridge Lemma
(\texttt{conj:bridge}, line~5466), upgrading averaged coherence to pointwise
coherence; and H2, a localised Constantin--Fefferman depletion criterion at
the shrinking scale \(\rstar\) --- no Type-I singularity can have
\(\liminf_{t\to T^-}\sstar(t)\le\varepsilon_0\). Every Type-I singularity must
maintain \emph{persistent geometric disorder}.

\begin{figure}[ht]
\centering
\begin{tikzpicture}[
  >={Latex[length=2mm]},
  early/.style={draw, rounded corners=2pt, font=\scriptsize, align=center,
                inner sep=3.5pt, text width=32mm, minimum height=9mm},
  surv/.style={early, draw=provedgreen, very thick},
  dead/.style={early, draw=impossiblered, dashed},
  new/.style={early, draw=openblue, thick, fill=panelgrey},
  ar/.style={->, thick}
]
\node[font=\small\bfseries] at (-3.0,2.6) {\Sn{29}--\Sn{30}};
\node[font=\small\bfseries] at (-3.0,-1.2) {\Sn{32}--\Sn{47}};

\node[dead] (rec)  at (0.0,3.3) {\S20 enstrophy recursion};
\node[dead] (thmE) at (4.0,3.3) {Theorem E: \(\sigma\to0\) at blow-up};
\node[dead] (aev)  at (8.0,3.3) {\texttt{lem:A-evol} coercivity};
\node[surv] (split) at (0.0,1.5) {\texttt{lem:grad-omega-split}};
\node[surv] (sig)   at (4.0,1.5) {\(\sigma=A_{\mathrm{loc}}/M^{3/2}\) invariance};
\node[surv] (d8)    at (8.0,1.5) {D8: \(\norm{\nabla u}_\infty\!\lesssim\!M\) is false};

\node[new] (mag)  at (0.0,-0.6) {\texttt{lem:direction-eq}\\ magnitude Caccioppoli};
\node[new] (sst)  at (4.0,-0.6) {\(\sstar\), \(\sstar\le16\sigma\)};
\node[new] (L)    at (8.0,-0.6) {\(\Lfac\): Wall 1};

\node[new] (pin)  at (2.0,-2.6) {the alignment pincer\\ \CONDITIONAL{H1, H2}};
\node[new] (open) at (6.6,-2.6) {\(\sstar\)-decay:\\ the single deep gap};

\draw[ar, provedgreen] (split)--(mag);
\draw[ar, provedgreen] (sig)--(sst);
\draw[ar, provedgreen] (d8)--(L);
\draw[ar] (mag)--(pin); \draw[ar] (sst)--(pin);
\draw[ar] (sst)--(open);
\draw[ar, impossiblered, dashed] (thmE) to[out=0,in=90,looseness=0.7]
  node[right, font=\scriptsize, pos=0.35, align=left, text=impossiblered]
  {its \emph{statement}\\ becomes the\\ open problem} (open.north);
\draw[ar, openblue, dashed] (L) -- (open.east);

\node[font=\scriptsize\itshape, text=rulegrey, align=center] at (4.0,-4.1)
  {solid green: certified by the audit and carried forward \quad
   dashed red: withdrawn};
\end{tikzpicture}
\caption[The inheritance graph]{What crossed the \Sn{31} boundary. Three
objects survived --- two theorems and one defect --- and each seeded a
distinct part of the later programme. The withdrawn Theorem~E crossed too, but
only as a question.}
\label{fig:inheritance}
\end{figure}

\subsection{The irony, stated plainly}
\label{sec:theorem-e-irony}

\begin{record}
Theorem~E (\texttt{thm:blowup-alignment}, line~4858) asserted, in \Sn{30},
that any blow-up forces \(\sigma(t)\to0\). Its proof rests on
\texttt{thm:sigma-gronwall}, which Defect~D5 shows is conditional on the
Beale--Kato--Majda criterion, and on Regime~I, which Defect~D6 shows requires
a viscosity larger than universal constants. It is not among the objects the
audit certifies as surviving.

\smallskip
The single deep open problem of the entire later programme --- named as such
in \Sn{38}, reframed as structural rigidity in \Sn{39}, and carried as
\texttt{target:disorder-depletion} (line~7490) --- is: \emph{force
\(\sstar(t)\to0\) along a putative Type-I blow-up sequence.}

\smallskip
These are the same statement. The programme's central unsolved problem is a
theorem it once believed it had proved, and whose proof it withdrew.
\end{record}

This is not a curiosity. It is the sharpest available measure of what the
\Sn{31} audit actually cost: not a step in an argument, but the argument's
conclusion, which then had to be rebuilt as a target and has resisted every
attempt since. Fifteen sessions of work \Sn{32}--\Sn{45} --- the level-set
Moser scheme, the magnitude Caccioppoli, the second rung, the
\(K\)-self-absorption identity, the localised Constantin--Fefferman
derivation, the Regime-A impossibility proof, the profile-rigidity reframe,
the annulus iteration, the band conjecture, four sessions of numerics ---
were all, directly or indirectly, attempts to recover it honestly.

\section{Second inheritance: Defect 8 becomes the Calder\'on--Zygmund
logarithm}
\label{sec:cz-inheritance}

This is the most important inheritance in the book, and the one the programme
came closest to losing.

\subsection{What D8 actually said}

The second half of Defect~D8 is a single sentence in
\texttt{sec:audit-s31}:

\begin{record}
``Moreover the stretching estimate invokes
\(\norm{\nabla u}_{L^\infty(B_{\rstar})}\le CM\), but Calder\'on--Zygmund
operators are unbounded on \(L^\infty\); the correct bound carries a
logarithmic factor, which matters precisely because \S22 relies on
log-corrected estimates.''
\end{record}

That is the whole of it. It is stated as a defect in one lemma. It is in fact
a statement about the architecture of every estimate the programme would go on
to write.

\subsection{Why it is structural, not local}

The programme's estimates all have the same shape. A budget is written for a
quantity built from the vorticity --- local enstrophy, angular energy, a
coupling density. The budget's dangerous term is vortex stretching, which
involves \(\nabla u\). And \(\nabla u\) is not available: only \(\omega\) is,
and \(\nabla u\) must be recovered from it by Biot--Savart. Biot--Savart is a
Calder\'on--Zygmund operator. Taking a supremum of a Calder\'on--Zygmund image
costs a logarithm --- the \(L^\infty\) endpoint fails --- and so
\[
  \norm{\nabla u}_{L^\infty} \;\lesssim\; M\,\Lfac,
  \qquad \Lfac = \log\bigl(e+\norm{u}_{H^3}/M\bigr).
\]
The factor is paid \emph{once}, at that step, and then inherited by everything
downstream. It is not a bookkeeping artefact of any single estimate, and no
amount of care in the downstream algebra removes it.

\subsection{The four appearances}

{\small
\begin{longtable}{@{}l p{0.32\textwidth} p{0.46\textwidth}@{}}
\caption[The life of the Calder\'on--Zygmund logarithm]{The obstruction's
career, from a clause in a defect list to the programme's central structural
barrier. Between \Sn{34} and \Sn{44} it was rediscovered three times as an
apparently new difficulty; \Sn{48} closed one of its three branches.}
\label{tab:cz-life}\\
\toprule
When & Where & What was said \\
\midrule
\endfirsthead
\toprule
When & Where & What was said \\
\midrule
\endhead
\Sn{31} & \texttt{sec:audit-s31}, D8
  & The bound \(\norm{\nabla u}_{L^\infty}\le CM\) is false; the correct
    bound carries a logarithm. \\[3pt]
\Sn{32} & \texttt{rem:marginality} (5443)
  & The connection is still live: the remark states that the logarithm ``is
    mandatory: Calder\'on--Zygmund operators are unbounded on \(L^\infty\) ---
    this was Defect~8 of the S31 audit''. The direction equation is
    consequently \emph{exactly critical} at the scale \(\rstar\), up to a
    logarithm. \\[3pt]
\Sn{33} & \texttt{lem:local-enstrophy-typeI} (5701), \(\Lfac\) defined at 5712
  & The factor enters the formal machinery and acquires its symbol. Every
    budget from here on carries it. \\[3pt]
\Sn{34} & \texttt{prop:crude-K} (5903)
  & \textbf{First rediscovery.} The coupling term re-enters at top order
    multiplied by \(\Lfac\): ``the crude route misses closing the loop by
    exactly one logarithm''. \\[3pt]
\Sn{41} & \texttt{rem:s41-shortfall} (8395)
  & \textbf{Second rediscovery.} Absorption of the annulus core requires
    \(C\Lfac\le\delta\nu\), which fails for fixed \(\nu\) as
    \(\Lfac\to\infty\); the shortfall is ``exactly the same single
    logarithm''. \\[3pt]
\Sn{44} & \texttt{prop:s44-linfty-route} (9209)
  & \textbf{Third rediscovery.} Both remaining conjectures close if and only
    if the correlation constant \emph{decays} like \(C_1/\Lfac\).
    Boundedness is not enough, because the hypothesis is used precisely to
    cancel the \(\Lfac\) the budget carries. \\[3pt]
\Sn{39} & \texttt{prop:transfer-audit}(d) (6941)
  & The one apparent escape: on a Type-I profile, ``the logarithmic factor
    \(\Lfac\) may be replaced by an absolute constant''. But this is
    \PROVEDMOD{(R)}, and hypothesis~(R) \emph{is} an assumed log-free gradient
    bound. The barrier is relocated, not removed. Wall~2.
    \emph{Amended at \Sn{48}: see the final row.} \\[3pt]
\Sn{45} & \texttt{scale\_diagnostics.py}
  & The decay \Sn{44} requires is measured. Slope of \(\log c_K\) on
    \(\log\Lfac\): \(-0.196\) (\(r^2=0.05\)) and \(+0.088\) (\(r^2=0.02\)),
    against the required \(-1\). Flat. \\[3pt]
\Sn{47} & Wall~1
  & Consolidated. Classification \OPEN{} and structural: no proof says the
    logarithm cannot be removed, but the three appearances establish that it
    is one obstacle seen three times, not three obstacles. And explicitly:
    the factor ``is not a later regression but the audit's finding carried
    forward''. \\[3pt]
\Sn{48} & \texttt{thm:s48-R} (\S\,\texttt{sec:hypR-s48})
  & \textbf{The escape is real after all --- on that branch only.}
    Hypothesis~(R) is proved, in a stronger form than it was stated:
    \(\|\nabla^k\partial_s^l u_\lambda\|_\infty\le
    C_{k,l}(c_I)|s|^{-(k+1)/2-l}\), uniformly in \(\lambda\). The route is
    Serrin's criterion in its \emph{strict} subcritical range --- the
    available bound is an \(L^\infty\) \emph{velocity} bound, and
    \(2/s'+3/s=0<1\), so boundedness suffices and smallness is never needed.
    Consequence: \texttt{prop:transfer-audit}(d) first clause \PROVED{};
    \textbf{Wall~2 \WITHDRAWN{}}. The profile route does not relocate
    \(\Lfac\) --- it never passes through the Calder\'on--Zygmund operator
    at all. Two cautions, both load-bearing: part~(e) inherits a
    pre-existing \emph{saturation} proviso and is only
    \textsc{proved mod saturation}; and Wall~1 itself is untouched, because
    \S33--\S37 need \(\nabla u\) measured against \(M(t)\), not
    \((T-t)^{-1}\), and those agree only on the saturated set. \\
\bottomrule
\end{longtable}}

\subsection{Why it was nearly lost}

\texttt{PROGRESS.md} recorded the \Sn{31} audit as having \emph{seven} grounds.
The omitted one was D8. For sixteen sessions, the programme's own navigational
summary --- the file every session was instructed to read first --- did not
contain the observation that its central obstruction had already been
identified.

The consequence was not a miscount. It was that each rediscovery looked like a
new difficulty. \Sn{34} reports missing closure ``by exactly one logarithm'' as
a finding about the crude \(K\)-route. \Sn{41} reports the annulus core falling
short by ``exactly the same single logarithm'' and notes, correctly, that this
is the known obstruction resurfacing --- but relates it to \Sn{34}, not to the
audit. \Sn{44} establishes that the closure needs decay rather than
boundedness, which is the sharpest form of the same statement.

Only in \Sn{47}, when the walls document was written under the rule that the
proof document governs, did the count get checked and the chain get closed.
Wall~0's final clause is the repair: ``Defect~D8 is the first appearance of
Wall~1 below; it was already flagged in the audit and is the reason the
logarithm is not a later regression.''

\begin{obsv}[The general lesson]
\label{obs:summary-loss}
A research programme's navigational summary is a lossy compression of its
record, and what it loses is not random: it loses the items that do not fit
the current narrative. In \Sn{31} the narrative was ``the recursion is broken'',
and seven of the eight defects are about the recursion. D8's second half is
about something else entirely --- an architectural feature of every future
estimate --- and it was the item dropped. The rule that the primary document
governs its own summaries is what recovered it, sixteen sessions late.
\end{obsv}

\section{Third inheritance: the withdrawal becomes the house rules}
\label{sec:withdrawal-to-rules}

The \Sn{31} audit is the direct ancestor of every rule in
Chapter~\ref{ch:house-rules}. The mapping is close enough to tabulate.

\begin{center}
\small
\begin{tabular}{@{}p{0.14\textwidth} p{0.30\textwidth} p{0.48\textwidth}@{}}
\toprule
Defect & Rule it produced & Mechanism \\
\midrule
D1--D3
  & Independent re-derivation before merge (\Sn{41})
  & Three arithmetic errors reached a public claim because nobody recomputed
    them. The gate requires the orchestrator to re-derive load-bearing steps
    by hand. It caught the inverted volume factor in \Sn{44}. \\[3pt]
D4
  & State the quantifier class and check every step against it
  & The \Sn{33} rebuild opens with a non-circularity remark: the pincer is a
    contradiction argument at a putative \emph{first} singularity, so all
    estimates run on the smooth interval and D4 cannot recur by design. \\[3pt]
D5
  & Do not assume an equivalent of the conclusion
  & Every statement in \Sn{32} onward is labelled \emph{Proved},
    \emph{Conditional} or \emph{Open target}, and conditional hypotheses are
    named. \\[3pt]
D8 (first half)
  & Never assert a coercivity without deriving it
  & The \Sn{35} ``second good term'' was asserted and then found by \Sn{36} to
    cancel exactly (\texttt{lem:w2-cancel}); Wall~7 records it as one of three
    instances of ``a favourable statement written down before the computation
    producing it was closed''. \\[3pt]
D8 (second half)
  & The primary document governs its summaries
  & \Sn{47}. See \S\ref{sec:cz-inheritance}. \\[3pt]
The whole episode
  & Report negatives as negatives
  & \Sn{41} returns a proof that its own assignment is impossible; \Sn{45}
    reports the measurement its own closure route needs as absent;
    \Sn{47} declines to upgrade Wall~1 from \OPEN{} to \IMPOSSIBLE{}. \\
\bottomrule
\end{tabular}
\end{center}

\section{What is still load-bearing, and what is dead}
\label{sec:load-bearing}

\begin{table}[ht]
\centering
\small
\begin{tabular}{@{}p{0.44\textwidth} p{0.44\textwidth}@{}}
\toprule
\textbf{Still load-bearing today} & \textbf{Dead} \\
\midrule
\texttt{lem:grad-omega-split}. Used in the \Sn{34} magnitude Caccioppoli and
the \Sn{36} Gram negativity; without it neither works.
  & \texttt{thm:uniform-enstrophy} and the two corollaries. Withdrawn, never
    repaired, and Wall~0 states what repair would require. \\[6pt]
The scale invariance of \(\sigma\), through
\(\sstar\le16\sigma\), and \(\sstar\) itself. Every session \Sn{32}--\Sn{45}
is about \(\sstar\).
  & \texttt{lem:A-evol}'s coercivity. Never derived; the later programme
    derives the direction equation from scratch instead. \\[6pt]
\texttt{thm:routeC}. Still the terminal step of any closure that produces a
\(\delta_0>0\).
  & \texttt{thm:sigma-gronwall} as a route to regularity. D5 is not
    repairable: the hypothesis is the conclusion. \\[6pt]
Defect~D8's second half, as \emph{the obstruction}. It is the thing every
open item in Parts~IV and VI is ultimately about.
  & Regime~A (bounded enstrophy near a singularity). \Sn{38} proves it
    impossible: at a genuine blow-up \(\limsup\Omega=\infty\). Not dormant ---
    \IMPOSSIBLE. \\[6pt]
The verification discipline, in full.
  & The \(c_K\) measurement programme as \emph{evidence}. Wall~6:
    \IMPOSSIBLE{} as a test. \\[6pt]
\textbf{Dormant, not dead:} \texttt{thm:shear-claimA},
\texttt{thm:near-shear}, \texttt{prop:theta-cacc},
\texttt{lem:pressure-traceless}, the Gagliardo--Nirenberg mean localisation,
the \(\Gamma(r)\) hierarchy, \texttt{lem:pressure-forcing}. All proved, none
refuted, none used since \Sn{31}. This is the largest body of correct
machinery in the document that nothing currently stands on.
  & \S19's gradient recursion. Superseded by a withdrawn argument, which
    leaves it in an odd position: it was replaced by something worse, and the
    replacement's failure does not restore it, because \S19's own exponent was
    sublinear. \\
\bottomrule
\end{tabular}
\caption[Load-bearing against dead]{The early work, sorted by whether anything
currently rests on it.}
\label{tab:load-bearing}
\end{table}

\begin{obsv}[The shape of the inheritance]
\label{obs:inheritance-shape}
Almost nothing that the early work \emph{claimed} survived. Almost everything
that the early work \emph{noticed} did. The split identity is an observation;
the scale invariance of \(\sigma\) is an observation; the failure of
Calder\'on--Zygmund on \(L^\infty\) is an observation. All three are still in
use. The theorems built on top of them --- the recursion, the induction, the
Gr\"onwall bound, Theorem~E --- are gone.

That asymmetry is the practical content of the \Sn{31} audit and the reason
this book keeps the early work rather than deleting it. A programme that
discarded its failed chapter would have discarded the three observations along
with the theorems, and the three observations are what the later programme is
made of.
\end{obsv}
